\documentstyle[%


righttag%


,11pt%


,amssymb%


,russian%               remove in Eng.


,izvru%                 Set izven% in Eng.


]{amsart}





% \APM = almost periodic measure. Seek'n'use if necessary, in \input'ed file as well


% {=} {<} {\subset} {\doteq} {\le} {\in}





\newcommand{\Hom}{\operatorname{Hom}}


\newcommand{\loc}{\operatorname{loc}}


\newcommand{\frm}{\operatorname{frm}}


\newcommand{\rpm}{\operatorname{rpm}}


\newcommand{\DIR}{\operatorname{DIR}}


\newcommand{\orb}{\operatorname{orb}}


\newcommand{\comp}{\operatorname{comp}}


\newcommand{\conv}{\operatorname{conv}}


\newcommand{\cl}{\operatorname{cl}}


\newcommand{\dist}{\operatorname{dist}}


\newcommand{\APM}{\operatorname{APM}}


\newcommand\esssup{\operatornamewithlimits{ess\,sup}}


\newcommand\limv{\mathop{\ov{\lim}}\limits}


\newcommand\limn{\mathop{\underline{\lim}}\limits}


\numberwithin{equation}{section}





\newcommand{\tts}[1]{}





\theoremstyle{plain}


\theorembodyfont{\it}


\newtheorem{thms}{\hskip\parindent Теорема}[section]


\newtheorem{lems}{\hskip\parindent Лемма}[section]





\theoremstyle{definition}


\theorembodyfont{\rm}


\newtheorem{cors}{\hskip\parindent Следствие}[section]


\newtheorem{defns}{\hskip\parindent Определение}[section]


\newtheorem{notes}{\hskip\parindent Замечание}[section]


\newtheorem{exams}{\hskip\parindent Пример}[section]








%\hoffset=-15mm%                  For Eng.


\setcounter{page}{29}


\god{2005}


\nomer{10~(521)} %                        Remove in Eng.


%\nomer{Vol.\,49, No.\,10}%               Set in Eng.


%\str{\pageref{begin}--\pageref{end}}%  Set in Eng.


\udk{517.977}





\author{\it А.Г.\,Иванов}


% NB:   ^^ remove in Eng.


\title{Динамическая система сдвигов и существование решения


задачи почти периодической оптимизации}


\thanks{}


%%% sooner, first two footnotes should B placed as thanks


%%% due to specificity of STY-file. Leave it 2 me








\date{}


%\pagestyle{myheadings}%         Set in Eng.


\pagestyle{plain} %              Remove in Eng.


\begin{document}


%\thispagestyle{plain}%          Set in Eng.


\maketitle


\label{begin}











В классе обобщенных управлений\footnote{О важности процедуры расширения или


овыпукления в задачах оптимального управления см., напр., [1]--[5], а


в игровых задачах --- [6], [7].}


(мер) приведены достаточные условия


существования оптимального процесса в смысле работы ([8], с.\,525) и в


терминах динамической системы сдвигов, определенной на множестве


допустимых процессов нелинейной системы управления, указаны достаточные


условия существования почти периодического (п.\,п.) магистрального


процесса, являющегося решением задачи оптимального


управления п.\,п. движениями.


\medskip


\setcounter{section}{1}


\setcounter{equation}{0}


\setcounter{lems}{0}


\setcounter{defns}{0}


\setcounter{notes}{0}





{\bf 1.} Пусть $\Bbb R^{n}$ --- $n$-мерное


евклидово пространство со стандартной нормой, $\Hom(\Bbb R^n)$


--- пространство линейных операторов


$L:\Bbb R^n\to\Bbb R^n$ с нормой


$|L|\doteq\max\limits_{|x|\le1}|Lx|$. Через


$\frak V_{n,\Bbb T}\doteq 


\frak V(\Bbb T\times\frak U,\Bbb R^n)$,


где $\frak U\in\comp(\Bbb R^m)$, а в качестве $\Bbb T$


рассматривается либо отрезок прямой, либо вся прямая, обозначим совокупность


функций $\varphi:\Bbb T\times\frak U\to\Bbb R^n$,


удовлетворяющих условиям


$\varphi(t,\cdot)\in C(\frak U,\Bbb R^n)$ при п.\,в.


$t\in\Bbb T$, для каждого $u\in\frak U$ отображение


$t\mapsto\varphi(t,u)$ измеримо (по Лебегу) и


существует такая функция


$\psi_{\varphi}\in L_1^{\loc}(\Bbb T,\Bbb R_+)$


($\Bbb R_+\doteq[0,\infty)$), что


$\max\limits_{u\in\frak U}|\varphi(t,u)|\le\psi_{\varphi}(t)$


при п.\,в. $t\in\Bbb T$.


В $\frak V_{n,\Bbb T}$ определена норма


$\|\cdot\|_{\frak V_{n,\Bbb T}}$ [2].


В дальнейшем $\frak V^{\loc}(\Bbb R


\times\frak U,\Bbb R^n)$ --- совокупность таких отображений


$\varphi:\Bbb R\times\frak U\to\Bbb R^n$, что


$\varphi\in\frak V_{n,\Bbb T}$ для каждого отрезка $\Bbb T\subset\Bbb R$,


и $(\frm(\frak U),|\cdot|_{w})$


--- нормированное пространство мер Радона на $\Bbb R^m$, носитель которых


содержится в $\frak U$ ([2], c.\,297), $\rpm(\frak U)$


--- его подмножество, состоящее из вероятностных мер Радона и


$\DIR(\frak U)$ --- совокупность мер Дирака $\delta_u$, сосредоточенных в


точках $u\in\frak U$; $\Bbb M(\Bbb T,\frm(\frak U))$ --- совокупность таких


измеримых отображений $\mu:\Bbb T\to(\frm(\frak U),|\cdot|_{w})$, что


$\|\mu\|_{\Bbb T}\doteq\esssup\limits_{t\in\Bbb T}


|\mu(t)(\frak U)|<\infty$


($|\mu(t)|(\frak U)$ --- вариация меры $\mu(t)\in\frm(\frak U)$),


$\cal M_{\Bbb T}\doteq \Bbb M(\Bbb T,\rpm(\frak U))$


и $\cal M_{\Bbb T}^{(1)}$ --- совокупность


измеримых отображений $t\mapsto\delta_{u(t)}\in\DIR(\frak U)


\subset(\frm(\frak U),|\cdot|_{w})$,


которая изоморфна $\cal U_{\Bbb T}$ --- множеству измеримых


отображений\footnote{Этот факт и


приведенные ниже свойства пространства $\Bbb M(\Bbb T,\frm(\frak U))$


в случае, когда в качестве $\Bbb T$


рассматривается отрезок прямой, приведены в [2] и без труда


переносятся на случай, когда $\Bbb T=\Bbb R$ (напр., [9]).}


$u:\Bbb T\to\frak U$ и, следовательно, каждое


$u(\cdot)\in\cal U_{\Bbb T}$ можно рассматривать как элемент из


$\cal M_{\Bbb T}^{(1)}$, отождествляя его с отображением


$t\mapsto\delta_{u(t)}\in\DIR(\frak U)$.


Поскольку $\frak V_{1,\Bbb T}^*\cong


\Bbb M(\Bbb T,\frm(\frak U))$,


то в $\Bbb M(\Bbb T,\frm(\frak U))$ возможно определить


(слабую) норму $\|\cdot\|_{w,\Bbb T}$, относительно которой множества


$\cal M_{\Bbb T}$ и $\frak S_1\doteq\{\mu\in\Bbb M


(\Bbb T,\frm(\frak U)):


\|\mu\|_{\Bbb T}\le 1\}$ компактны, причем, если


$\mu_j\in\frak S_1$, $j\in\Bbb Z_+$, то


$\lim\limits_{j\to\infty}\|\mu_j-\mu_0\|_{w,\Bbb T}=0$ в том и только


том случае, если


$\lim\limits_{j\to\infty}\int\limits_{\Bbb T}


\langle\mu_j(t),\varphi(t,u)\rangle dt=


\int\limits_{\Bbb T}\langle\mu_0(t),\varphi(t,u)\rangle dt$


$\Big(\langle\mu_j(t),\varphi(t,u)\rangle\doteq


\int\limits_{\frak U}\varphi(t,u)\mu_j(du)$,


$j\in\Bbb Z_{+}\Big)$ для каждой функции $\varphi\in\frak V_{1,\Bbb T}$.





Сейчас фиксируем функцию $f: G\times\frak U\to\Bbb R^n$ ($G$


--- область в $\Bbb R^n$), удовлетворяющую условию





1)~в каждой точке $(x,u)\in G\times\frak U$ существует производная


$f'_x(x,u)$ и $f\in C(G\times\frak U,\Bbb R^n)$,


$f'_x\in C(G\times\frak U,\Hom(\Bbb R^n))$, и рассмотрим систему уравнений


\begin{equation}


\Dot x=\langle\mu(t),f(x,u)\rangle dt


\doteq\int_{\frak U}f(x,u)\mu(t)(du),\ \


\mu(\cdot)\in\cal M_{\Bbb T},


\end{equation}


для которой через $\frak A_c(\Bbb T)$


обозначим совокупность (допустимых) пар $(x(\cdot),\mu(\cdot))$, в


которых $x(t)$, $t\in\Bbb T$, --- решение системы (1.3), отвечающее


(обобщенному) управлению $\mu(\cdot)\in\cal M_{\Bbb T}$, такое, что


$\ov\orb(x;\Bbb T)\subset G$, где $\ov\orb(x;\Bbb T)$ --- замыкание


в $\Bbb R^n$ множества $\orb(x;\Bbb T)


\doteq\{x(t),\ t\in\Bbb T\}$ и для заданного


$X\in\comp(\Bbb R^n)$ полагаем $\frak A_c(\Bbb T;X)\doteq


\{(x(\cdot),\mu(\cdot))\in\frak A_c(\Bbb T):


\ov\orb(x;\Bbb T)\subset X\}$.





Из определения $\frak A_c(\Bbb T;X)$,


используя теорему Арцеля--Асколи, компактность множества


$\cal M_{\Bbb T}\subset(\Bbb M(\Bbb T,\frm(\frak U)),


\|\cdot\|_{w,\Bbb T})$, учитывая, что равенство


$\lim\limits_{j\to\infty}\|\mu_j - \mu_0\|_{w,\Bbb T}=0$,


$\{\mu_j\}_{j\in\Bbb Z_+}\subset\cal M_{\Bbb T}$, где


$\Bbb T\doteq[t_0,t_1]$, влечет (напр., [3]) равенство


$\lim\limits_{j\to\infty}\Big(\max\limits_{t\in[t_0,t_1]}\int\limits_{t_0}^t


\langle\mu_j(t)-\mu_0(t),\varphi(t,u)\rangle dt\Big)=0$ для


каждой функции $\varphi\in\frak V_{n,\Bbb T}$,


получаем следующее утверждение.








\begin{lems} 


Для каждого отрезка $\Bbb T\subset\Bbb R$


множество $\frak A_c(\Bbb T;X)


\subset C(\Bbb T,\Bbb R^n)\times\cal M_{\Bbb T}$


$($сек\-вен\-ци\-аль\-но\/$)$ компактно.


\end{lems}





Далее, на $\frak A_c[t_0,t_1]$ зададим функционал


\begin{equation}


(x(\cdot),\mu(\cdot))\mapsto


\frak T(x(\cdot),\mu(\cdot);t_0,t_1)


\doteq\int_{t_0}^{t_1}\langle\mu(t),g(x(t),u)\rangle dt,


\end{equation}


где $g\in C(G\times\frak U,\Bbb R)$, который, как


несложно показать, непрерывен на $C(\Bbb T,\Bbb R^n)


\times\cal M_{\Bbb T}$. Поэтому по лемме 1.1 в силу


теоремы Вейерштрасса получается








\begin{lems}


Для каждого отрезка $[t_0,t_1]$ $($глобальное\/$)$ решение задачи


\begin{equation}


\frak T(x(\cdot),\mu(\cdot);t_0,t_1)\to\min,\ \ (x(\cdot),\mu(\cdot))


\in\frak A_c([t_0,t_1];X)


\end{equation}


существует.


\end{lems}





Совокупность решений задачи (1.3) обозначим через $\cal R[t_0,t_1]$.





По аналогии с определением в ([8], с.\,24) дадим





\begin{defns}


Допустимый процесс


$(\wh x(\cdot),\wh \mu(\cdot))\in\frak A_c[t_0,t_1]$


системы (1.1) называется оптимальным для задачи


$\frak T(x(\cdot),\mu(\cdot);t_0,t_1)\to\min,\ (x(\cdot),\mu(\cdot))


\in\frak A_c[t_0,t_1]$, если для любого другого процесса


$(x(\cdot),\mu(\cdot))\in\frak A_c[t_0,t_1]$ такого, что


$x(t_0)=\wh x(t_0)$, $x(t_1)=\wh x(t_1)$, выполнено неравенство


$\frak T(x(\cdot),\mu(\cdot);t_0,t_1)


\ge\frak T(\wh x(\cdot),\wh\mu(\cdot);t_0,t_1)$.


\end{defns}





Совокупность решений указанной в определении~1.1 задачи


обозначим через $O\cal P[t_0,t_1]$ и положим


\begin{equation} 


O\cal P(\Bbb R)\doteq


\{(\wh x(\cdot),\wh\mu(\cdot))\in\frak A_c(\Bbb R):


\text{ для каждого } [t_0,t_1]\subset\Bbb R


\ (\wh x(\cdot),\wh\mu(\cdot))|_{[t_0,t_1]}


\in O\cal P[t_0,t_1]\}.


\end{equation}


Совокупность решений задачи~(1.3) в смысле определения~1.1 обозначим


через $O{\cal P}([t_0,t_1];X)$ и аналогично $O{\cal P}({\Bbb R})$


определяем множество $O{\cal P}({\Bbb R};X)


\subset{\frak A}_c(\Bbb R;X)$.








\begin{notes}


Поскольку $\cal R[t_0,t_1]\subset O\cal P([t_0,t_1];X)$, а сужение


$(x(\cdot),\mu(\cdot))|_{[t'_0,t'_1]}$ решения


$(x(\cdot),\mu(\cdot))\in O\cal P([t_0,t_1];X)$ на каждый подотрезок


$[t'_0,t'_1]\subset[t_0,t_1]$ принадлежит


$O\cal P([t'_0,t'_1];X)$ ([8], c.\,24), то и сужение пары


$(x(\cdot),\mu(\cdot))\in\cal R[t_0,t_1]$


на каждый подотрезок $[t'_0,t'_1]\subset[t_0,t_1]$ также принадлежит


$O\cal P([t'_0,t'_1];X)$.


\end{notes}





Всюду далее, не оговаривая, предполагаем, что $X\in\comp(G)$ такое, что


для каждого отрезка $\Bbb T\subset\Bbb R$ множество


$\frak A_c(\Bbb T;X)\ne\emptyset$.


Такое множество существует, если, например, система (1.1) имеет процесс


$(x(\cdot),\mu(\cdot))\in\frak A_{c}(\Bbb R)$, в котором


$\ov\orb(x)\doteq\ov\orb(x;\Bbb R)


\in\comp(G)$. В этом случае в качестве множества $X$ достаточно


рассмотреть любую компактную окрестность $\ov\orb(x)$,


содержащуюся в~$G$.








\begin{lems}


Для любой системы отрезков $\Bbb T_1\subset\Bbb T_2\subset\dotsb$,


исчерпывающей $\Bbb R$, существует


последовательность $(x_j(\cdot),\mu_j(\cdot))


\in O\cal P(\Bbb T_j;X)$, $j\in\Bbb N$, сходящаяся на каждом


отрезке $\Bbb T\subset\Bbb R$ в


топологии пространства $C(\Bbb T,\Bbb R^n)


\times\cal M_{\Bbb T}$ к некоторому процессу


$(\wh x(\cdot),\wh\mu(\cdot))\in\frak A_c(\Bbb R;X)$.


\end{lems}





\begin{pf}


Рассмотрим последовательность


$\{(\wh x_j(\cdot),\wh\mu_j(\cdot))\}_{j=1}^\infty$,


cоставленную из решений (см. лемму 1.2) задачи (1.3) на


$\frak A_c(\Bbb T_j;X)$,


$j\in\Bbb N$. По лемме 1.1 из этой последовательности можно извлечь


подпоследовательность


$\{(\wh x{}_j^{(1)}(\cdot),\wh\mu{}_j^{(1)}(\cdot))\}_{j=1}^\infty$,


сходящуюся при $j\to\infty$ в пространстве $C(\Bbb T_1)\times\cal M_{\Bbb T_1}$


$(C(\Bbb T_1)\doteq C(\Bbb T_1,\Bbb R^n))$


к некоторому процессу $(z_1(\cdot),\nu_1(\cdot))\in\frak A_c(\Bbb T_1;X)$


и, кроме того (см. замечание 1.1),


$(\wh x{}_j^{(1)}(\cdot),\wh\mu{}_j^{(1)}(\cdot))\in O\cal P(\Bbb T_1;X)$,


$j\in\Bbb N$. Снова по лемме 1.1 из последовательности


$\{(\wh x{}_j^{(1)}(\cdot),\wh\mu{}_j^{(1)}(\cdot))\}_{j=1}^\infty$


извлекаем подпоследовательность


$(\wh x{}_j^{(2)}(\cdot),\wh\mu{}_j^{(2)}(\cdot))\in O\cal P(\Bbb T_2;X)$,


$j\in\Bbb N$, сходящуюся при $j\to\infty$ в пространстве


$C(\Bbb T_2)\times\cal M_{\Bbb T_2}$


к некоторому процессу $(z_2(\cdot),\nu_2(\cdot))\in\frak A_c(\Bbb T_2;X)$.


Продолжив указанную процедуру, получим совокупность последовательностей


$\{(\wh x{}_j^{(i+1)}(\cdot),\wh\mu{}_j^{(i+1)}(\cdot))\}_{j=1}^\infty


\subset\{(\wh x{}_j^{(i)}(\cdot),\wh\mu{}_j^{(i)}(\cdot))\}_{j=1}^\infty$,


$i\in\Bbb N$, такую, что


$(\wh x{}_j^{(i)}(\cdot),\wh\mu{}_j^{(i)}(\cdot))\in


O\cal P(\Bbb T;X)$, $j\in\Bbb N$,


при каждом $i\in\Bbb N$ и всяком $\Bbb T\subset\Bbb T_i$


и сходящуюся при $j\to\infty$


в пространстве $C(\Bbb T_i)\times\cal M_{\Bbb T_i}$


к некоторому процессу $(z_i(\cdot),\nu_i(\cdot))\in\frak A_c(\Bbb T_i;X)$.


Поскольку $(z_i(\cdot),\nu_i(\cdot))|_{\Bbb T_l}=


(z_l(\cdot),\nu_l(\cdot))$, $l=1,\dotsc,i-1$, то процесс


$(\wh x(\cdot),\wh\mu(\cdot))\in\frak A_c(\Bbb R;X)$ такой, что


$(\wh x(\cdot),\wh\mu(\cdot))|_{\Bbb T_k}=


(z_k(\cdot),\nu_k(\cdot))$, $k\in\Bbb N$, и последовательность


$(x_j(\cdot),\mu_j(\cdot))\doteq(\wh x_j^{(j)}(\cdot),


\wh\mu_j^{(j)}(\cdot))$, $j\in\Bbb N$, искомые.


\end{pf}





\addtocounter{section}{1}


\setcounter{equation}{0}


\setcounter{footnote}{0}


\setcounter{lems}{0}


\setcounter{defns}{0}


\setcounter{notes}{0}


\setcounter{thms}{0}





{\bf 2.} Далее используем следующие два определения.








\begin{defns}


Пусть заданы процесс


$(\wh x(\cdot),\wh\mu(\cdot))\in\frak A_c(\Bbb R)$ и


числа $\varepsilon$, $\vartheta$, $\eta>0$. Система (1.1) называется


$\varepsilon, \vartheta, \eta$-равномерно


локально управляемой ($\varepsilon,\vartheta,\eta$-РЛУ) на


$\orb(\wh x;[\tau_0,\infty))$, если при каждом $\tau\ge\tau_0$ для любого


$x_0\in O_{\varepsilon}[\wh x(\tau)]


\doteq\{x\in\Bbb R^n:|\wh x(\tau)-x|\le\varepsilon\}$


найдется такое управление $\mu\in\cal M_{\tau,\vartheta}


\doteq\cal M_{[\tau,\tau+\vartheta]}$, что


$\|\wh\mu-\mu\|_{w,[\tau,\tau+\vartheta]}


\le\eta|\wh x(\tau)-x_0|$ и система (1.1) имеет решение


$x(t)$, $\tau\le t\le\tau+\vartheta$, удовлетворяющее


условиям $x(\tau)=x_0$, $x(\tau+\vartheta)=\wh


x(\tau+\vartheta)$. Эта система называется РЛУ на $\orb(\wh


x;[\tau_0,\infty))$, если существуют такие константы $\varepsilon$,


$\vartheta$, $\eta>0$, что она является


$\varepsilon,\vartheta,\eta$-РЛУ на $\orb(\wh x;[\tau_0,\infty))$.


\end{defns}








\begin{defns}


Пусть заданы процесс


$(\wh x(\cdot),\wh\mu(\cdot))\in\frak A_c(\Bbb R)$ и


числа $\varepsilon$, $\vartheta$, $\eta>0$. Система (1.1) называется


$\varepsilon,\vartheta,\eta$-равномерно локально достижимой


($\varepsilon,\vartheta,\eta$-РЛД)


с


$\orb(\wh x;[\tau_0,\infty))$,


если при каждом $\tau\ge\tau_0$ для любого


$x_0\in O_{\varepsilon}[\wh x(\tau+\vartheta)]$ найдется такое управление


$\mu\in\cal M_{\tau,\vartheta}$, что


$\|\wh\mu-\mu\|_{w,[\tau,\tau+\vartheta]}


\le\eta|\wh x(\tau+\vartheta)-x_0|$


и система (1.1) имеет решение $x(t)$,


$\tau\le t\le\tau+\vartheta$, удовлетворяющее условиям


$x(\tau)=\wh x(\tau)$, $x(\tau+\vartheta)=x_0$.


Система (1.1) называется РЛД с $\orb(\wh x;[\tau_0,\infty))$,


если существуют такие константы $\varepsilon$, $\vartheta$, $\eta>0$, что она


является $\varepsilon,\vartheta,\eta$-РЛД


с $\orb(\wh x;[\tau_0,\infty))$.


\end{defns}








\begin{notes}


Поскольку при каждом $s\in\Bbb R$


функция $t\mapsto\wh x_{s}(t)\doteq\wh x(t+s)$, $t\in[0,\vartheta]$,


является решением задачи Коши


\begin{gather*}


\Dot x=\langle\wh\mu_{s}(t),f(x,u)\rangle,


\ \ \wh\mu_{s}(t)\doteq\wh\mu(t+s),\ \ t\in[0,\vartheta];\\


x(0)=\wh x_{s}(0), 


\end{gather*}


то система (1.1) \ $\varepsilon,\vartheta,\eta$-РЛУ


на $\orb(\wh x;[\tau_0,\infty))$ в том и только том случае,


если при каждом $\tau\ge\tau_0$ для любого $x_0


\in O_{\varepsilon}[\wh x(\tau)]$


найдется такое $\nu\in\cal M_{[0,\vartheta]}$, что


$\|\wh\mu_\tau-\nu\|_{w,[0,\vartheta]}


\le\eta|\wh x(\tau)-x_0|$


и система (1.1) имеет решение $y(t)$,


$0\le t\le\vartheta$, удовлетворяющее условиям


$y(0)=x_0$, $y(\vartheta)=\wh x(\tau+\vartheta)$.





Аналогичное утверждение имеет место и для РЛД системы (1.1) с


$\orb(\wh x;[\tau_0,\infty))$.


\end{notes}





В дальнейшем важную роль будет играть свойство одновременной


РЛУ системы (1.1) на $\orb(\wh x;[\tau_0,\infty))$ и ее РЛД с


$\orb(\wh x;[\tau_0,\infty))$.


Для формулировки достаточных условий выполнения этого свойства


рассмотрим систему


\begin{equation}


\Dot x=A(t)x+\langle\Delta\mu(t),f(\wh x(t),u)\rangle,


\quad A(t)\doteq\langle\wh\mu(t),f'_{x}(\wh x(t),u)\rangle,\ \


\Delta\mu(t)\doteq\wh\mu(t)-\mu(t),


\end{equation}


которая называется РЛУ (в нуль) на $[\tau_0,\infty)$, если найдутся такие


константы $\varepsilon$, $\vartheta>0$, что для каждого


${\tau\ge\tau_0}$ и всякого $x_0\in O_{\varepsilon}[0]$


найдется такое $\mu\in\cal M_{\tau,\vartheta}$, при котором


система (2.1) имеет решение $x(t)$, $\tau\le


t\le\tau+\vartheta$, удовлетворяющее условиям


$x(\tau)=x_0$, $x(\tau+\vartheta)=0$.





Для фиксированной пары $(\wh x(\cdot),\wh\mu(\cdot))\in


\frak A_c(\Bbb R;X)$ выбираем такое $r>0$,


что компактное множество


\begin{equation}


K\doteq\ov\orb(\wh x;\Bbb R)+O_r[0]\subset G.


\end{equation}





Для непрерывной вместе со своей производной по $x$


функции $f:G\times\frak U\to\Bbb R^n$ предполагаем выполненным условие





2)~существуют такие константы $\wh


r\in(0,r]$, $\alpha>0$ и непрерывная ограниченная функция


$(t,u)\mapsto\frak f(t,u)\in[\tau_0,\infty)$, что для всех


$(t,z,u)\in[\tau_0,\infty)\times O_{\Hat r}[0]\times\frak U$


выполнено неравенство


$\max\limits_{\theta\in [0,1]}|f'_x(\wh x(t)+\theta


z,u)-f'_x(\wh x(t),u)|\le\frak f(t,u)|z|^\alpha$.





В точности повторив доказательства теорем 2.1 и 4.1 работы


[10], заменив $\orb_+(\wh x)


\doteq\orb(\wh x;[0,\infty))$ на


$\orb(\wh x;[\tau_0,\infty))$, с учетом сделанного в ней замечания 4.2


получаем следующее утверждение.








\begin{thms}


Пусть функция $f:G\times\frak U\to\Bbb R^n$ удовлетворяет


условию~$1)$ и на процессе $(\wh x(\cdot),\wh\mu(\cdot))


\in\frak A_c(\Bbb R;X)$


выполнено условие~$2)$. Тогда, если система $(2.1)$ РЛУ на $[\tau_0,\infty)$,


то найдутся такие константы $\varepsilon$, $\vartheta$, $\eta>0$, что система


$(1.1)$ \ $\varepsilon,\vartheta,\eta$-РЛУ на


$\orb(\wh x;[\tau_0,\infty))$


и


$\varepsilon,\vartheta,\eta$-РЛД


с


$\orb(\wh x;[\tau_0,\infty))$. При этом для каждого


$\tau\ge\tau_0$ решение, фигурирующее в определении


$\varepsilon,\vartheta,\eta$-РЛУ системы $(1.1)$ на


$\orb(\wh x;[\tau_0,\infty))$


и


$\varepsilon,\vartheta,\eta$-РЛД с $\orb(\wh x;[\tau_0,\infty))$, при всех


$t\in[\tau,\tau+\vartheta]$ содержится в $K$.


\end{thms}





В дальнейшем используем следующий вариант леммы~2.2 из [10].








\begin{lems}


Пусть выполнены условия теоремы $2.1$ и последовательность


$\{x_0^{(j)}\}_{j=1}^\infty\subset O_\varepsilon[\wh x(\tau)]$


$(\tau\ge\tau_0)$ такая, что


$\lim\limits_{j\to\infty}|x_0^{(j)}-\wh x(\tau)|=0$.


Пусть далее $\mu_j\in\cal M_{\tau,\vartheta}$ --- управление,


удовлетворяющее неравенству


$\|\wh\mu-\mu_j\|_{w,[\tau,\tau+\vartheta]}\le\eta|\wh


x(\tau)-x_0^{(j)}|$,


при котором система $(1.1)$ имеет такое решение $x_j(t)\in K$, что


$x_j(\tau)=x_0^{(j)}$, $x_{j}(\tau+\vartheta)=\wh


x(\tau+\vartheta)$. Тогда


$\lim\limits_{j\to\infty}\|x_j-\wh x\|_{C[\tau,\tau+\vartheta]}=0$.


\end{lems}





Аналогично доказательству упомянутой леммы из [10] доказывается








\begin{lems}


Пусть выполнены условия теоремы $2.1$ и последовательность


$\{x_0^{(j)}\}_{j=1}^\infty


\subset O_\varepsilon[\wh x(\tau+\vartheta)]$


$(\tau\ge\tau_0)$ такая, что


$\lim\limits_{j\to\infty}|x_0^{(j)}-\wh x(\tau+\vartheta)|=0$.


Пусть далее $\mu_j\in\cal M_{\tau,\vartheta}$ --- управление,


удовлетворяющее неравенству


$\|\wh\mu-\mu_j\|_{w,[\tau,\tau+\vartheta]}\le\eta|\wh


x(\tau+\vartheta)-x_0^{(j)}|$,


при котором система $(1.1)$ имеет такое решение $x_j(t)\in K$, что


$x_j(\tau)=\wh x(\tau)$, $x_j(\tau+\vartheta)=x_0^{(j)}$. Тогда


$\lim\limits_{j\to\infty}\|x_j-\wh x\|_{C[\tau,\tau+\vartheta]}=0$.


\end{lems}








\begin{thms}


Пусть функция $f:G\times\frak U\to\Bbb R^n$ удовлетворяет


условию~$1)$ и последовательность


$\{(x_j(\cdot),\mu_j(\cdot))\}_{j=1}^\infty$,


$(x_j(\cdot),\mu_j(\cdot))\in O\cal P(\Bbb T_{j})$,


$\Bbb T_{j}\subset\Bbb T_{j+1}$, $j\in\Bbb N$


$\big(\mathop{\cup}\limits_{j=1}^{\infty}\Bbb T_{j}=\Bbb R\big)$,


сходится на каждом отрезке $\Bbb T\subset\Bbb R$ в пространстве


$C(\Bbb T)\times\cal M_{\Bbb T}$ к процессу


$(\wh x(\cdot),\wh\mu(\cdot))


\in\frak A_c(\Bbb R)$, на котором выполнено условие~$2)$.


Тогда, если система $(2.1)$ РЛУ на $[\tau_0,\infty)$, то


найдется такое $T_1>\tau_0$, что $(\wh


x(\cdot),\wh\mu(\cdot))$ принадлежит $O\cal P[T_1,\infty)$.


\end{thms}








\begin{pf}


По теореме 2.1 найдутся такие константанты


$\varepsilon$, $\vartheta$, $\eta>0$, что система (1.1) будет


$\varepsilon,\vartheta,\eta$-РЛУ на $\orb(\wh


x;[\tau_0,\infty))$


и


$\varepsilon,\vartheta,\eta$-РЛД


с


$\orb(\wh x;[\tau_0,\infty))$. Полагаем $T_1\doteq\tau_0+\vartheta$ и


покажем, что $(\wh x(\cdot),\wh\mu(\cdot))


\in O\cal P[T_1,\infty)$. Действительно, в противном случае


найдется отрезок $[t_0,t_1]\subset[T_1,\infty)$ и такой процесс


$(x(\cdot),\mu(\cdot))\in\frak A_c[t_0,t_1]$, что


\begin{equation}


x(t_0)=\wh x(t_0),\ x(t_1)=\wh x(t_1).


\end{equation} 


При некотором $\delta>0$ выполнено неравенство (см. обозначение в (1.2))


\begin{equation}


\frak T(x(\cdot),\mu(\cdot);t_0,t_1)<


\frak T(\wh x(\cdot),\wh\mu(\cdot);t_0,t_1)-\delta.


\end{equation}





В силу условий теоремы 2.2 будет выполняться равенство


\begin{equation}


\lim_{j\to\infty}(\|\wh x-x_{j}\|_{C(\Bbb T)}+


\|\wh\mu-\mu_j\|_{w,\Bbb T})=0,\ \ \Bbb T\doteq[t_0-\vartheta,t_1+\vartheta].


\end{equation}


Поэтому для всех $j$, начиная с некоторого $\wh j\in\Bbb N$,


$x_j(t)\in O_{\varepsilon}[\wh x(t)]$ для всех $t\in\Bbb T$.


Фиксируем теперь $j\ge\wh j$ и для точки $x_j(t_0-\vartheta)


\in O_\varepsilon[\wh x(t_0-\vartheta)]$


($t_0-\vartheta\ge\tau_0$ в силу выбора $T_1$)


рассмотрим $\mu_{j}^{+}\in\cal M_{[t_0-\vartheta,t_0]}$, удовлетворяющее


неравенству


\begin{equation}


\|\wh\mu-\mu_{j}^+\|_{w,[t_0-\vartheta,t_0]}\le\eta


|\wh x(t_0-\vartheta)-x_j(t_0-\vartheta)|,


\end{equation}


и система (1.1) имеет решение $y_j^+(t)\in O_\varepsilon[\wh x(t)]$,


$t\in[t_0-\vartheta,t_0]$, с условиями


\begin{equation}


y_j^+(t_0-\vartheta)=x_j(t_0-\vartheta),\ \ y_j^+(t_0)=\wh x(t_0).


\end{equation}


Для точки $x_j(t_1+\vartheta)\in O_\varepsilon[\wh x(t_1+\vartheta)]$


рассмотрим $\mu_j^-\in\cal M_{[t_1,t_1+\vartheta]}$, удовлетворяющее


неравенству


\begin{equation}


\|\wh\mu-\mu_j^-\|_{w,[t_1,t_1+\vartheta]}\le\eta


|\wh x(t_1+\vartheta)-x_j(t_1+\vartheta)|,


\end{equation}


и система (1.1) имеет решение $y_j^-(t)\in O_\varepsilon[\wh x(t)]$,


$t\in[t_0,t_1+\vartheta]$, с условиями


\begin{equation}


y_j^-(t_1)=\wh x(t_1),\ \ y_j^-(t_1+\vartheta)=x_j(t_1+\vartheta).


\end{equation}


Теперь определим 


$$


\nu_j(t)\doteq


\begin{cases}


\mu_j^+(t), & t\in[t_0-\vartheta,t_0]; \\


\mu(t), & t\in[t_0,t_1]; \\


\mu_j^-(t), & t\in[t_1,t_1+\vartheta], \ \nu_j\in\cal M_{\Bbb T}.


\end{cases}


$$


Тогда (см. (2.3), (2.7) и (2.9)) отвечающее ему решение


$y_j^+(t)$, $t\in\Bbb T$, системы (1.1) имеет вид


$$


y_j(t)\doteq


\begin{cases}


y_{j}^{+}(t), & t\in[t_0-\vartheta,t_0]; \\


y(t), & t\in[t_0,t_1]; \\


y_j^-(t), & t\in[t_1,t_1+\vartheta],


\end{cases}


$$


причем


\begin{equation*}


y_j(t_0-\vartheta)=x_j(t_0-\vartheta),


\ \ y_j(t_1+\vartheta)=x_j(t_1+\vartheta).


\end{equation*}





Далее, в силу определения процесса $(y_j(\cdot),\nu_j(\cdot))$ имеем


соотношения


\begin{multline*}


\frak T(y_j(\cdot),\nu_j(\cdot);t_0-\vartheta,t_1+\vartheta)=


\frak T(y_j^+(\cdot),\mu_j^+(\cdot);t_0-\vartheta,t_0)+


\frak T(x(\cdot),\mu(\cdot);t_0,t_1)+ \\


%


+\frak T(y_j^-(\cdot),\nu_j^-(\cdot);t_1,t_1+\vartheta)\overset{(2.4)}{<}{}


\frak T(\wh x(\cdot),\wh


\mu(\cdot);t_0-\vartheta,t_1+\vartheta)-\delta+I_{j}^{(1)}+I_{j}^{(2)},


\end{multline*}


в которых


\begin{align*} 


I_j^{(1)}&\doteq\int_{t_0-\vartheta}^{t_0}(\langle\mu_j^+(t),


g(y_j^+(t),u)\rangle-\langle\wh\mu(t),g(\wh x(t),u)\rangle)dt,\\


I_j^{(2)}&\doteq\int_{t_1}^{t_1+\vartheta}(\langle\mu_j^-(t),


g(y_j^-(t),u)\rangle-\langle\wh\mu(t),g(\wh x(t),u)\rangle)dt.


\end{align*}





Покажем, что


\begin{equation}


\lim_{j\to\infty}I_j^{(k)}=0,\ \ k=1,2.


\end{equation}


Действительно,


\begin{equation}


|I_j^{(1)}|\le\vartheta\omega_{\gamma_j}[g,K\times\frak U]+


\bigg|\int_{t_0-\vartheta}^{t_0}


\langle\mu_j^+(t)-\wh\mu(t),g(\wh x(t),u)\rangle dt\bigg|,


\end{equation}


где $\gamma_j\doteq\|y_j^+-\wh x\|_{C[t_0-\vartheta,t_0]}$,


а


$\omega_{\gamma_j}[g,K\times\frak U]$ --- $\gamma_j$-колебание


непрерывной функции $g$ на компакте $K\times\frak U$.


По лемме 2.1 в силу (2.5) и (2.7) \ $\gamma_j\to0$


при $j\to\infty$. Следовательно,


\begin{equation}


\lim_{j\to\infty}\omega_{\gamma_j}[g,K\times\frak U]=0,


\end{equation}


а т.\,к.


$\lim\limits_{j\to\infty}\|\mu_j^+


-\wh\mu\|_{w,[t_0-\vartheta,t_0]}=0$ (см. (2.5) и (2.6)), то


$\lim\limits_{j\to\infty}\Big|\int\limits_{t_0-\vartheta}^{t_0}


\langle\mu_j^+(t)-\wh\mu(t),g(\wh x(t),u)\rangle dt\Big|=0$.


Из последнего равенства, учитывая (2.11) и (2.12), получаем (2.10) для


$k=1$. Аналогично, используя лемму 2.2, равенства (2.5), (2.9) и


неравенство (2.8), доказываем (2.10) при $k=2$.


Поэтому найдется такое $j_1\ge\wh j$, что при всех $j\ge j_1$ будет


выполнено неравенство


\begin{equation}


\frak T(y_j(\cdot),\nu_j(\cdot);t_0-\vartheta,t_1+\vartheta)<


\frak T(\wh x(\cdot),\wh \mu(\cdot);t_0-\vartheta,t_1+\vartheta)-\delta/2.


\end{equation}


С другой стороны, 


в силу (2.5) \ $\lim\limits_{j\to\infty}\frak T(x_j(\cdot),


\mu_j(\cdot);t_0-\vartheta,t_1+\vartheta)=


\frak T(\wh x(\cdot),\wh \mu(\cdot);t_0-\vartheta,t_1+\vartheta)$.


Следовательно,


$\frak T(\wh x(\cdot),\wh \mu(\cdot);t_0-\vartheta,t_1+\vartheta)<


\frak T(x_{j_0}(\cdot),\mu_{j_0}(\cdot);t_0-\vartheta,t_1+\vartheta)+\delta/4$


при некотором $j_0\ge j_1$. Отсюда, учитывая (2.13), получаем


$$


\frak T(y_{j_0}(\cdot),\nu_{j_0}(\cdot);t_0-\vartheta,t_1+\vartheta)<


\frak T(x_{j_0}(\cdot),\mu_{j_0}(\cdot);t_0-\vartheta,t_1+\vartheta)


-\delta/4<\frak T(x_{j_0}(\cdot),\mu_{j_0}(\cdot);t_0-\vartheta,t_1+\vartheta),


$$


что противоречит тому, что $(x_{j_0}(\cdot),\mu_{j_0}(\cdot))\in


O\cal P[t_0-\vartheta,t_1+\vartheta]$.


\end{pf}








\begin{notes}


Как показано в лемме~1.3, расширение множества $\cal U_{\Bbb T}$


до $\cal M_{\Bbb T}$ позволяет по


любой системе отрезков $\Bbb T_1\subset\Bbb T_2\subset\dotsb$,


исчерпывающей $\Bbb R$, указать


последовательность $(x_j(\cdot),\mu_j(\cdot))


\in O{\cal P}({\Bbb T}_j;K)$, $j\in{\Bbb N}$, сходящуюся на


каждом отрезке ${\Bbb T}\subset{\Bbb R}$ в топологии пространства


$C(\Bbb T,\Bbb R^n) \times\cal M_{\Bbb T}$ к


некоторому процессу $(\wh x(\cdot),\wh\mu(\cdot))\in


\frak A_c(\Bbb R;K)$. Как видно из приведенного


доказательства теоремы~2.2, этот процесс будет принадлежать множеству


$O\cal P([T_1,\infty);K)$, если указанные при доказательстве этой


теоремы решения $y_j^+$ и $y_j^-$, помимо приведенных свойств,


еще не должны покидать множество $K$, соответственно, при всех


$t\in[t_0-\vartheta,t_0]$ и $t\in[t_1,t_1+\vartheta]$. Проверка такого


свойства и достаточные условия его выполнения требуют, вообще говоря,


дополнительного исследования. Поэтому в рамках статьи ограничимся


рассмотрением оптимального процесса, принадлежащего


$O{\cal P}[T_1,\infty)$, указанного в теореме~2.2.


\end{notes}





В дальнейшем, при описании поведения оптимальных процессов,


определенных на $[0,T]$, при $T\to\infty$ (что является одной из основных


задач теории магистральных процессов [8] и асимптотически


достижимых элементов [11]) важную роль будет играть динамическая система сдвигов,


определенная на множестве $\frak A_c(\Bbb R;X)$. 


\medskip


\addtocounter{section}{1}


\setcounter{equation}{0}


\setcounter{lems}{0}


\setcounter{defns}{0}


\setcounter{notes}{0}


\setcounter{thms}{0}


\setcounter{cors}{0}


\setcounter{exams}{0}





{\bf 3.}


Пусть $(\frak Y,\rho)$ --- метрическое пространство.


Через $L_1^{\loc}(\Bbb R,\frak Y)$ обозначим совокупность таких


измеримых отображений $f:\Bbb R\to\frak Y$, что при некотором (а


значит, и при любом) $y\in\frak Y$ функция $t\mapsto\rho(y,f(t))$


принадлежит пространству $L_1^{\loc}(\Bbb R,\Bbb R)$. На


$L_1^{\loc}(\Bbb R,\frak Y)$ введем расстояние


\begin{equation}


\varrho(f,g)\doteq


\sup_{t\in\Bbb R}\bigg[\min\bigg\{\int_t^{t+1}


\rho(f(s),g(s))ds,\ |t|^{-1}\bigg\}\bigg].


\end{equation}


Несложно показать, что пространство


\begin{equation}


\frak B\doteq(L_1^{\loc}(\Bbb R,\frak Y),\varrho)


\end{equation}


является метрическим и 


неравенство $\varrho(f,g)<\varepsilon$ равносильно выполнению неравенства\\


$\max\limits_{|t|\le\varepsilon^{-1}}


\int\limits_t^{t+1}\rho(f(s),g(s))ds<\varepsilon$.


Отсюда в силу неравенств


$$


\max_{|t|\le\vartheta}


\int_t^{t+1}\rho(f(s),g(s))ds\le


\frak p_{\vartheta}(f,g)\ \ (\vartheta>0),


\quad \frak p_{\vartheta}(f,g)


\le(2\vartheta+1)\max_{|t|\le\vartheta}


\int_t^{t+1}\rho(f(s),g(s))ds\ \ (\vartheta\in\Bbb N),


$$


где здесь и всюду далее


\begin{equation}


\frak p_{\vartheta}(f,g)\doteq


\int_{-\vartheta}^{\vartheta+1}\rho(f(s),g(s))ds


\quad (\vartheta>0),


\end{equation}


вытекает








\begin{lems}


Пусть $(\Bbb A,\prec)$ --- направленное множество,


функции $f$ и $f_{\alpha}$, $\alpha\in\Bbb A$, принадлежат


$L_1^{\loc}(\Bbb R,\frak Y)$.


Тогда $\lim\limits_{\alpha\in\Bbb A}\varrho(f_{\alpha},f)=0$ в том и


только том случае, если для каждого $\vartheta>0$ выполнено равенство


$\frak p_{\vartheta}(f_{\alpha},f)=0$


или, что равносильно, $\lim\limits_{\alpha\in\Bbb A}


\Big(\max\limits_{|t|\le\vartheta}\int\limits_t^{t+1}


\rho(f_{\alpha}(s),f(s))ds\Big)=0$.


\end{lems}





С использованием леммы 3.1, теоремы Фреше [12] об изометрии


сепарабельного метрического пространства $\frak Y$ части некоторого


сепарабельного банахового пространства $\Bbb B$, которая является замкнутой,


если $\frak Y$ --- полное сепарабельное метрическое пространство, и


свойств нормированного пространства


$(L_1(\Bbb T,\Bbb B),\|\cdot\|_{L_1})$,


где $\Bbb T$ --- числовой промежуток, доказывается








\begin{thms}


Пусть метрическое пространство $\frak Y$ сепарабельно. Тогда


$\lim\limits_{\tau\to0}


\varrho(f(\cdot+{\tau}),f(\cdot)){=}0$


для каждой функции


$f\in L_1^{\loc}(\Bbb R,\frak Y)$ 


и если к тому же пространство $\frak Y$ полное,


то метрическое пространство $\frak B$ также является полным.


\end{thms}








\begin{notes}


Отметим, что на $L_1^{\loc}(\Bbb R,\frak Y)$


определено также $d$-расстояние (или иначе\linebreak $d_{\rho}$-расстояние)


\begin{equation}


d(f,g)\doteq\sup_{t\in\Bbb R}\int_t^{t+1}\rho(f(s),g(s))ds,


\quad f,g\in L_1^{\loc}(\Bbb R,\frak Y).


\end{equation}


\end{notes}





Функция


$f\in L_1^{\loc}(\Bbb R,\frak Y)$ называется $d$-ограниченной, если


$d(\varphi,y){<}\infty$ при некотором $y\in\frak Y$. По аналогии со случаем


$\frak Y=\Bbb R$ [13] совокупность


$\Bbb S(\Bbb R,\frak Y)$ \ $d$-ограниченных функций из


$L_1^{\loc}(\Bbb R,\frak Y)$ называем пространством Степанова.


Поскольку $\varrho(f,g)\le d(f,g)$ для любых


$f,g\in L_1^{\loc}(\Bbb R,\frak Y)$, то


топология $\cal T_d$ на $\Bbb S(\Bbb R,\frak Y)$,


индуцированная метрикой $d$, содержится в топологии


$\cal T_\varrho$, индуцированной метрикой $\varrho$.


Отметим также, если $(\frak Y,\rho)$ --- полное сепарабельное метрическое


пространство, то пространство


$(\Bbb S(\Bbb R,\frak Y),d)$ является полным.





Далее, на метрическом пространстве $\frak B$ (см. (3.2))


введем однопараметрическое семейство отображений


$g^{t}:\frak B\to\frak B$,


определенных равенством $g^t(f)=f_t$, $f\in\frak B$, где


$f_t(\cdot)\doteq f(\cdot+t)$ --- сдвиг на $t\in\Bbb R$ функции


$f(\cdot)$, и в дальнейшем через $\orb_g(f)$ и


$\orb_g^+(f)$ обозначаем соответственно траекторию и положительную


полутраекторию заданного движения $t\mapsto g^{t}(f)$.





Из леммы 3.1 и теоремы 3.1 вытекает








\begin{thms}


$(\frak B,g^{t})$ --- динамическая система $($в смысле Маркова\/$)$.


\end{thms}








\begin{notes}


Введенная на $L_1^{\loc}(\Bbb R,\frak B)$ метрика $\varrho$ (см. (3.1))


является естественным распространением метрики Бебутова $\varrho_c$


[14], заданной на пространстве непрерывных функций $C(\Bbb R,\frak Y)$


равенством


$\varrho_{c}(f,g)\doteq\sup\limits_{t\in\Bbb R}


[\min\{\rho(f(t),g(t)),\frac{1}{|t|}\}]$, $f,g\in C(\Bbb R,\frak Y)$,


в которой неравенство $\varrho_c(f,g)<\varepsilon$ равносильно тому, что


$\max\limits_{|t|\le\varepsilon^{-1}}\rho(f(t),g(t))<\varepsilon$ и,


следовательно, $\lim\limits_{j\to\infty}\varrho_c(f,f_{j})=0$


в том и только том случае, если для любых $\varepsilon$, $\vartheta>0$


найдется такое $j_0=j_0(\varepsilon,\vartheta)\in\Bbb N$, что


$\max\limits_{|t|\le\vartheta}\rho(f_j(t),f(t))<\varepsilon$


при $j\ge j_0$. Теорема 3.2 утверждает, что так же, как и в


непрерывном случае, пара $(\frak B,g^t)$ является динамической


системой (динамической системой сдвигов). Отметим далее, что


динамическая система сдвигов на пространстве непрерывных функций,


описанная, например, в [14]--[16], играет важную роль в теории


динамических систем и ее приложениях (см., напр., [14]--[19] и


приведенную там библиографию). В теории управляемых систем она


по-видимому впервые стала использоваться в [20] в связи с


введенными Е.Л.\,Тонковым понятиями равномерной локальной и глобальной


управляемости, равномерной наблюдаемости и стабилизации линейных систем


управления. Динамическая система сдвигов на множестве функций с


введенной на нем топологией посредством оператора замыкания


использовалась в [8], [21] при изучении магистральных процессов.


\end{notes}





В следующей теореме и далее, если не оговорено специально,


рассматриваем метрическое пространство $\frak B$, определеное равенством


(3.2), и $\cl A$ --- замыкание множества $A\subset\frak B$.








\begin{thms}


Пусть $(\frak Y,\rho)$ --- полное сепарабельное метрическое пространство и


$\cal K{\subset}\frak B$. Тогда $\cl\cal K\in


\comp(\frak B)$ в том и только том случае, если множество $\cal K$


равностепенно \label{foot3} не\-пре\-рыв\-но\footnote{То есть при каждом


$\vartheta>0$ (см. обозначение (3.3)) $\lim\limits_{\tau\to


0}(\sup\limits_{f\in\cal K}\frak p_\vartheta


(f(\cdot+\tau),f(\cdot))ds)=0$.}


и существует такое $y\in\frak Y$, что


$\sup\limits_{f\in\cal K}


\frak p_\vartheta(f,y)ds<\infty$ для любого $\vartheta>0$.


\end{thms}


\setcounter{footnote}{0}








\begin{pf}


Необходимость условий теоремы 3.3 является


очевидным следствием теоремы Хаусдорфа о предкомпактности


подмножества метрического пространства [12] и теоремы~3.1.





Для доказательства достаточности надо использовать теорему


Фреше [12] об изометрии полного сепарабельного


метрического пространства $(\frak Y,\rho)$ и замкнутого множества


некоторого сепарабельного банахова пространства


$\Bbb B$, а затем для расширяющейся системы отрезков


$[-\vartheta,\vartheta+1]$, $\vartheta\in\Bbb N$,


последовательно применить теорему Рисса ([12], с.\,40) о


предкомпактности множества в


$(L_1([-\vartheta,\vartheta+1],\Bbb B),\|\cdot\|_{L_1})$.


\end{pf}





Отметим далее (см. замечание 3.1),


что из неравенства $\frak p_\vartheta(f,y)\le


2([\vartheta]+1)d(f,y)$ ($\vartheta>0$,\linebreak $y\in\frak Y$) вытекает, что


для всякого $d$-ограниченного множества функций $F\subset


L_1^{\loc}(\Bbb R,\frak Y)$, т.\,е. такого, что


$\sup\limits_{\varphi\in F}d(\varphi,y)\doteq\frak d<\infty$


при некотором $y\in\frak Y$,


для каждого $\vartheta>0$ выполняется неравенство


$\sup\limits_{\varphi\in F}\frak p_\vartheta(\varphi,y)


\le2([\vartheta]+1)\frak d$. Поэтому из


теоремы 3.3 получаем








\begin{cors}


Если равностепенно непрерывное множество $\cal K\subset\frak B$


является $d$-огра\-ни\-чен\-ным, то $\cl\cal K\in\comp(\frak B)$.


\end{cors}





Напомним (см. замечание 3.1), что $\Bbb S(\Bbb R,\frak Y)\subset\frak B$.








\begin{thms}


Пусть $f\in\Bbb S(\Bbb R,\frak Y)$. Тогда множество


$\cl(\orb_g(f))\subset\frak B$ является $d$-огра\-ни\-чен\-ным и принадлежит 


$\comp(\frak B)$ в том и только том случае, если $f$ \ $d$-непрерывна.


\end{thms}








\begin{pf}


Пусть $\cal K{\doteq}\cl(\orb_g(f))$ (напомним, что


$\orb_g(f){\doteq}\{g^{t}(f)$, $t\in\Bbb R\}$ и $g^t(f){\doteq}


f_t$). По определению для каждой функции $\wh f\in\cal K$


найдется такая числовая последовательность $\{t_j\}_{j=1}^\infty$, что


\begin{equation}


\lim_{j\to\infty}\varrho(g^{t_j}(f),\wh f)=0.


\end{equation}


Поэтому для каждого $t\in\Bbb R$ существует такое $j=j(t)\in\Bbb N$, что


(см. обозначение (3.3)) $\frak p_{|t|}(f_{t_j},\wh f)


\le 1$ и, следовательно, из неравенства


$\int\limits_t^{t+1}\rho(\wh f(s),y)ds\le


\frak p_{|t|}(f_{t_j},\wh f)+d(f,y)$


($y\in\frak Y$), вытекает, что $d(\wh f,y)\le 1+d(f,y)$,


т.\,е. множество $\cal K$ \ $d$-ограничено.





Допустим теперь, что $f$ \ $d$-непрерывна. По определению для любого


$\varepsilon>0$ найдется такое $\delta\in(0,1]$, что


$d(f_{\tau},f)\le\varepsilon/2$ при


всех $\tau\in(-\delta,\delta)$. Фиксируем далее произвольную функцию


$\wh f\in\cal K$. В силу (3.5) по лемме 3.1 для каждого


$\vartheta=|t|+1$, $t\in\Bbb R$, найдется такое $t_j$, что


$\frak p_\vartheta(g^{t_j}(f),\wh f)<\varepsilon/4$. Поэтому из соотношений


$\int\limits_t^{t+1}\rho(\wh f_{\tau}(s),\wh f(s))ds\le


\frak p_\vartheta(g^{t_j}(f),\wh f)+d(f_{\tau},f)<\varepsilon$


получаем, что $d(\wh f_{\tau},\wh f){\le}\varepsilon$,


т.\,е. множество $\cal K$ \ $d$-равностепенно непрерывно, а значит 


(см. замечание 3.1), и равностепенно непрерывно в метрике $\varrho$. 


По следствию 3.1 \ $\cal K\in\comp(\frak B)$.





Обратно, если $\cal K\in\comp(\frak B)$, то по 


теореме 3.3 \ $\cal K$ равностепенно непрерывно. Поэтому 


(см. сноску на этой же странице и лемму 3.1)


%%%%%%% ^^^^^ refer to as Footnote~3


для любого $\varepsilon>0$ найдется такое $\delta>0$, что при всех


$\tau\in(-\delta,\delta)$ и


каждой функции $\wh f\in\cal K$ \ $\frak p_1(\wh f_{\tau},\wh f)<\varepsilon$.


Учитывая, что $g^t(f)\in\cal K$ при всех $t\in\Bbb R$,


получаем (см. (3.3) при $\vartheta=1$), что


$d(f_{\tau},f)\le\varepsilon$ при $|\tau|<\delta$, т.\,е. $f$


является $d$-непрерывной.


\end{pf}


\setcounter{footnote}{0}





В дальнейшем понадобятся рекуррентные и


п.\,п. движения определенной


динамической системы $(\frak B,g^t)$. Считаем далее,


что в определении пространства $\frak B$ (см. (3.2))


$(\frak Y,\rho)$ --- полное сепарабельное метрическое пространство.








\begin{defns}


Функция $f\in\frak B$ называется рекуррентной, если для любого


$\varepsilon>0$ множество


$\{\tau\in\Bbb R:\varrho(f_{\tau},f)<\varepsilon\}$


относительно плотно (на $\Bbb R$).


\end{defns}





Совокупность рекуррентных функций в $\frak B$ обозначим через


$\frak R(\Bbb R,\frak Y)$.





Из определения метрики $\varrho$ и леммы 3.1 вытекает, что


$f\in\frak R(\Bbb R,\frak Y)$ в том и только том случае,


если для любых $\varepsilon,\vartheta>0$ (см. (3.3)) множество


$E_{\frak R}(f,\varepsilon,\vartheta)


\doteq\{\tau\in\Bbb R:\frak p_\vartheta(f_\tau,f)<\varepsilon\}$


ее $\varepsilon,\vartheta$-почти периодов ($\varepsilon,\vartheta$-п.\,п.)


относительно плотно.





Приведем необходимые в дальнейшем утверждения о структуре множества


$\frak R(\Bbb R,\frak Y)$.








\begin{lems}


Пусть $f\in\frak R(\Bbb R,\frak Y)$. Тогда


$\lim\limits_{h\to0}d(f_h,f)=0$


и при каждом фиксированном $y\in\frak Y$ $($см. $(3.4))$ $d(f,y)<\infty$.


\end{lems}








\begin{pf}


Фиксируем произвольное $t\in\Bbb R$ и пусть $l>0$ ---


константа, входящая в определение относительной плотности множества


$E_{\frak R}(f,1,|t|)$. Тогда при каждом\linebreak


$\tau\in E_{\cal R}(f,1,|t|)\cap [-t,-t+l]$


имеем следующие соотношения:


$$


\int_t^{t+1}\rho(f(s),y)ds\le


\frak p_{|t|}(f_{\tau},f)+


\int_t^{t+1}\rho(f_{\tau}(s),y)ds


<1+\int_0^{l+1}\rho(f(s),y)ds\doteq{\frak k}


$$


и, следовательно, $d(f,y)\le\frak k$.





Далее, для заданного $\varepsilon>0$ возьмем точку


$\tau\in E_{\frak R}(f,\frac{\varepsilon}{2},-|t|-1)\cap [-t,-t+l]$,


где $l>0$ --- число, входящее в определение относительной плотности


множества $E_{\frak R}(f,\frac{\varepsilon}{2},-|t|-1)$, и выберем


(см. теорему 3.1) $\delta\in(0,1]$ таким, что


$\frak p_{l+1}(f_{h},f)<\varepsilon/4$ при $|h|\le\delta$. 


Тогда при $|h|\le\delta$ из неравенств


$\int\limits_t^{t+1}\rho(f_{h}(s),f(s))ds\le


2\frak p_{|t|+1}(f_{\tau},f)+


\frak p_{l+1}(f_{h},f)<2\frac{\varepsilon}{4}+\frac{\varepsilon}{2}


=\varepsilon$


получаем, что $d(f_{h},f)\le\varepsilon$ при $|h|\le\delta$.


\end{pf}





Из доказанной леммы 3.2 и теоремы 3.4 (см. также ее доказательство)


получаем








\begin{cors}


Множество $\frak R(\Bbb R,\frak Y)\subset\Bbb S(\Bbb R,\frak Y)$ и


для каждой функции $f\in\frak R(\Bbb R,\frak Y)$


мно\-же\-ство\footnote{Напомним, что


множество $\Bbb S(\Bbb R,\frak Y)\subset L_1^{\loc}(\Bbb R,\frak Y)$


определено в замечании 3.1 и далее рассматриваем его как подпространство


пространства $\frak B$.}


$\cl(\orb_g(f))\in\comp(\frak B)$ и $d$-равностепенно непрерывно.


\end{cors}


\setcounter{footnote}{0}





Кроме того, в силу этой же


леммы~3.2 получается








\begin{lems}


Метрическое пространство $(\frak R(\Bbb R,\frak Y),d)$ является полным.


\end{lems}








\begin{pf}


Так как метрическое пространство (см. замечание 3.1)


$(\Bbb S(\Bbb R,\frak Y),d)$ является полным, то для


доказательства леммы 3.3 достаточно показать, что если


последовательность $\{f_j\}_{j=1}^\infty\subset


\frak R(\Bbb R,\frak Y)$ такова, что


$\lim\limits_{j\to\infty}d(f_{j},f)=0$, то $f\in\frak R(\Bbb R,\frak Y)$.


Действительно, если $\lim\limits_{j\to\infty}d(f_{j},f)=0$, то,


используя неравенство $\frak p_{\vartheta}(g^\tau(f),f)\le


2d(f_j,f)+\frak p_\vartheta(g^\tau(f_j),f_{j})$, выполненное для


всех $j\in\Bbb N$, $\tau\in\Bbb R$ и $\vartheta>0$, получим, что


для любых $\varepsilon,\vartheta>0$ и


каждом достаточно большом $j$ относительно плотное множество


$E_{\frak R}(f_j,\frac{\varepsilon}{2},\vartheta)$ содержится в


$E_{\cal R}(f,\varepsilon,\vartheta)$.


\end{pf}





По теореме 3.2 \ $(\frak B,g^t)$ --- динамическая


система, поэтому можно определить рекуррентные и п.\,п. движения


[15]. Следуя [15], говорим, что движение


$t\mapsto g^{t}(f)$, отвечающее $f\in\frak B$, рекуррентно, если для любого


$\varepsilon>0$ найдется такое $l>0$, что для каждого $t\in\Bbb R$ при


любом $t_0\in\Bbb R$ \ $\{\tau\in\Bbb R:\varrho(g^{t+\tau}(f),g^{t}(f))<


\varepsilon\}\cap[t_0,t_0+l]\ne\emptyset$.





В силу следствия 3.2 и теоремы 2.9 из ([15], с.\,405) вытекает








\begin{lems}


Функция $f\in\frak R(\Bbb R,\frak Y)$


в том и только том случае, если движение $t\mapsto g^{t}(f)$ рекуррентно.


\end{lems}





Напомним [15], что движение $t\mapsto g^{t}(f)$ называется


п.\,п., если для любого $\varepsilon>0$ множество


$\{\tau\in\Bbb R:\sup\limits_{t\in\Bbb R}


\varrho(g^{t+\tau}(f),g^t(f))<\varepsilon\}$


относительно плотно и функция $f\in L_1^{\loc}(\Bbb R,\frak Y)$ называется


п.\,п. по Степанову [13], если для каждого $\varepsilon>0$ множество


$E_{S}(f,\varepsilon)\doteq\{\tau\in\Bbb R:d(f_{\tau},f)<\varepsilon\}$


ее $\varepsilon$-п.\,п. относительно плотно.





Совокупность п.\,п. по Степанову функций из


$L_1^{\loc}(\Bbb R,\frak Y)$ обозначим $S(\Bbb R,\frak Y)$. Так же, как в


скалярном случае [13], можно показать, что всякая функция из


$S(\Bbb R,\frak Y)$ является $d$-ограниченной. Поэтому


(см. замечание 3.1) $S(\Bbb R,\frak Y)\subset\Bbb S(\Bbb R,\frak Y)$.








\begin{lems}


Движение $t\mapsto g^t(f)\in\frak B$ п.\,п. в том и только том случае,


если $f\in S(\Bbb R,\frak Y)$.


\end{lems}








\begin{pf}


Пусть движение $t\mapsto g^{t}(f)$ п.\,п. и $\tau$ такое, что


$\sup\limits_{t\in\Bbb R}\varrho(g^{t+\tau}(f),g^{t}(f))<\varepsilon$.


Поэтому (см. (3.1))


$\int\limits_t^{t+1}\rho(f_{\tau}(s),f(s))ds\le


\max\limits_{|\xi|\le\varepsilon^{-1}}


\int\limits_\xi^{\xi+1}\rho(g^{t+\tau}(f)(s),g^t(f)(s))ds<\varepsilon$


при каждом $t\in\Bbb R$.


Следовательно, каждый $\varepsilon$-п.\,п. движения $t\mapsto g^{t}(f)$


принадлежит множеству $E_{S}(f,\varepsilon)$, а значит,


$f\in S(\Bbb R,\frak Y)$.





Обратно, если $f\in S(\Bbb R,\frak Y)$, то из неравенства


$\sup\limits_{t\in\Bbb R}\varrho(g^{t+\tau}(f),g^t(f))


\le d(f_{\tau},f)$,


выполненного для любого фиксированного $\tau\in\Bbb R$ вытекает, что


движение $t\mapsto g^{t}(f)$ п.\,п.


\end{pf}








\begin{lems}


Пусть функция $f\in S(\Bbb R,\frak Y)$. Тогда замыкание


$\cal H(f)$ множества $\orb_g(f)$ в метрике $d$ совпадает с


$\cl(\orb_g(f))$.


\end{lems}








\begin{pf}


Поскольку (см. замечание~3.1) $\cal T_d\subset\cal T_\varrho$, то


$\cal H(f)\subset\cl(\orb_g(f))$. Докажем, что верно и


обратное включение. Действительно, по лемме 3.4 движение


$t\mapsto g^t(f)$ п.\,п. Следовательно [15], движение


$t\mapsto g^t(\wh f)$, отвечающее функции $\wh f\in\cl(\orb_g(f))$,


является п.\,п. и поэтому $\wh f\in S(\Bbb R,\frak Y)$.


Так как $f, \wh f\in S(\Bbb R,\frak Y)$, то для любого


$\varepsilon>0$ существует такое $l>0$,


что при каждом $t\in\Bbb R$ найдется


$\tau\in E_S(f,\varepsilon/3)\cap E_S(\wh f,\varepsilon/3)\cap[-t,-t+1]$.


Далее, т.\,к. $\wh f\in\cl(\orb_g(f))$, то найдется такая


последовательность $\{t_{j}\}_{j=1}^\infty\subset\Bbb R$, что


$\lim\limits_{j\to\infty}\varrho(g^{t_j}(f),\wh f)=0$.


Поэтому по лемме 3.1 найдется такое $j_0\in\Bbb N$


(см. обозначение (3.3)), что для всех


$j\ge j_0$ будет выполнено неравенство


$\frak p_l(g^{t_j}(f),\wh f)<\varepsilon/3$.


Но тогда при каждом $t\in\Bbb R$, учитывая выбор


$\tau$, для всех $j\ge j_0$ получаем, что


$\int\limits_t^{t+1}\rho(g^{t_j}(f)(s),\wh f(s))ds\le


d(g^\tau(\wh f),\wh f)+d(g^\tau(f),f)+


\frak p_l(g^{t_j}(f),\wh f)<\varepsilon$.


Отсюда следует, что 


$d(g^{t_j}(f),\wh f)\le\varepsilon$ при всех $j\ge j_0$,


т.\,е. $\wh f\in\cal H(f)$. Таким образом, доказано, что


$\cl(\orb_g(f))\subset\cal H(f)$ и что совместно с обратным включением


$\cal H(f)\subset\cl(\orb_g(f))$


завершает доказательство леммы~3.6.


\end{pf}





Приведем два примера при конкретном выборе пространства $\frak Y$.








\begin{exams} 


На множестве $\Phi\doteq\{(A,V)\in\Hom(\Bbb R^n)


\times\conv(\Bbb R^n):0\in V\}$ определим


расстояние $\rho_\Phi(\varphi_1,\varphi_2)\doteq|A_1-A_2|


+\dist(V_1,V_2)$, $\varphi_k=(A_k,V_k)\in\Phi$, $k=1,2$, где


$\dist$ --- метрика Хаусдорфа


на $\conv(\Bbb R^n)$. Тогда $(\Phi,\rho_{\Phi})$ --- полное


сепарабельное метрическое пространство. Рассмотрим множество


$\Bbb S(\Bbb R,\Phi)\subset L_1^{\loc}(\Bbb R,\Phi)$


(см. замечание 3.1). Как и в [22], [23], каждую


функцию $t\mapsto\varphi(t)=(A(t),V(t))$ из $\Bbb S(\Bbb R,\Phi)$


отождествим с системой управления


$\Dot x=A(t)x+v$, $(x,v)\in\Bbb R^n\times V(t)$, в


которой допустимыми управлениями служат измеримые сечения отображения


$t\mapsto V(t)$, содержащего при п.\,в. $t\in\Bbb R$ нуль и при этом (см.


определение множества $\Phi$) не исключается возможность того, что нуль


при некоторых или п.\,в. $t$ принадлежит границе $V(t)$.


(О целесообразности изучения такого класса систем см., напр.,


[10], [22], [24] и приведенную там библиографию.) Часть результатов


данного раздела была анонсирована в [23], доказана в [25], когда в


качестве $\frak B$ рассматривалось пространство


$(\Bbb S(\Bbb R,\Phi),\varrho)$, и


свойства динамической системы сдвигов, используемые в указанных работах,


доказывались с использованием структуры пространства


$\Bbb S(\Bbb R,\Phi)$. Отметим, что в силу теоремы 3.1 требование


$\lim\limits_{\tau\to 0}\rho(\varphi_{\tau},\varphi)=0$ для получения


результатов, указанных в ([23], c.\,467), излишне.


\end{exams}








\begin{exams}


Рассмотрим случай, когда в качестве полного


сепарабельного метрического пространства $(\frak Y,\rho)$ берется (см.


раздел~1) $(\frm(\frak U),\rho_w)$, где $\rho_w$ --- метрика,


индуцированная нормой $|\cdot|_w$ на $\frm(\frak U)$.


Отметим, что $\Bbb M(\Bbb R,\frm(\frak U))\subset\Bbb S(\Bbb R,\frm(\frak U))$


и (см. (3.1) при $\frak Y=\frm(\frak U)$) на


$\Bbb M(\Bbb R,\frm(\frak U))$ определено расстояние


$$


\varrho_w(\mu,\nu)\doteq\sup_{t\in\Bbb R}


\bigg[\min\bigg\{\int_t^{t+1}


\rho_w(\mu(s),\nu(s))ds,\ |t|^{-1}\bigg\}\bigg],


\ \ \mu, \nu\in\Bbb M(\Bbb R,\frm(\frak U)),


$$


в котором (см. лемму 3.1) равенство


$\lim\limits_{j\to\infty}\varrho_w(\mu,\mu_j)=0$ равносильно тому, что


$$


\lim_{j\to\infty}


\int_{-\vartheta}^{\vartheta+1}\rho_w(\mu(s),\mu_j(s))ds=0


\ \ \text{при каждом $\vartheta>0$}.


$$


Поэтому с использованием определения нормы $\|\cdot\|_{w,\Bbb T}$ на


$\Bbb M(\Bbb T,\frm(\frak U))$ несложно доказывается


\end{exams}








\begin{lems}


Пусть $\mu,\mu_j\in\cal M_{\Bbb R}$,


$j\in\Bbb N$, и $\lim\limits_{j\to\infty}\varrho_w(\mu,\mu_j)=0$.


Тогда при каждом $\vartheta>0$


$$


\lim\limits_{j\to\infty}


\|\mu-\mu_{j}\|_{w,[-\vartheta,\vartheta+1]}=0.


$$


\end{lems}





Отметим (см. замечание 3.1), что на


$\Bbb M(\Bbb R,\frm(\frak U))$ определено также $d_{w}$-расстояние


\begin{equation}


d_w(\mu,\nu)\doteq\sup_{t\in\Bbb R}\int_t^{t+1}


\rho_w(\mu(s),\nu(s))ds,\ \ \mu,\nu\in\Bbb M(\Bbb R,\frm(\frak U)).


\end{equation}


\medskip   %%%%% <-- can B removed in Eng.


\addtocounter{section}{1}


\setcounter{equation}{0}


\setcounter{lems}{0}


\setcounter{defns}{0}


\setcounter{notes}{0}


\setcounter{thms}{0}


\setcounter{cors}{0}


\setcounter{exams}{0}





{\bf 4.} В этом разделе и далее рассматриваем метрическое пространство


\begin{equation}


\frak P\doteq(C_{b}(\Bbb R,


\Bbb R^n)\times\cal M_{\Bbb R},\varrho_{с,w}),


\end{equation}


где $C_b(\Bbb R,\Bbb R^n)$ --- совокупность ограниченных функций из


$C(\Bbb R,\Bbb R^n)$ с метрикой $\varrho_{с,w}$, индуцированной (см.


замечание 3.1 и пример 3.2) метриками $\varrho_с$ и $\varrho_w$,


определенными на $C_b(\Bbb R,\Bbb R^n)$ и подмножестве


$\cal M_{\Bbb R}\subset\Bbb M(\Bbb R,\frm(\frak U))$


соответственно. Так как $\Bbb R^n$ и


$\cal M_{\Bbb R}$ --- полные сепарабельные


метрические пространства, то (см. [16] и теорему 3.1 при


$\frak Y\doteq\rpm(\frak U)$) пространство $\frak P$,


определенное в (4.1), также является полным сепарабельным


метрическим пространством, и пара


(см. [16] и теорему 3.2 при $\frak Y\doteq\rpm(\frak U)$)


$(\frak P,g^t)$, где $g^t(x(\cdot),\mu(\cdot))


\doteq(x_t(\cdot),\mu_t(\cdot))$,


$(x(\cdot),\mu(\cdot))\in\frak P$, является динамической системой.


Отметим, что множество $\frak A_c(\Bbb R,X)


\subset C_b(\Bbb R,\Bbb R^n)\times\cal M_{\Bbb R}$


инвариантно (относительно потока $g^t$).





Всюду далее $\cl P$ --- замыкание множества $P\subset\frak P$


и в следующих двух леммах $\varepsilon,\vartheta,\eta$-РЛУ в малом


на $\orb(\wh x;[\tau_0,\infty))$


и $\varepsilon,\vartheta,\eta$-РЛД в малом


с $\orb(\wh x;[\tau_0,\infty))$ означают, что


выполнено свойство решений, указанное во второй части теоремы 2.1.








\begin{lems}


Пусть $q(\cdot)\doteq(\wh x(\cdot),


\wh\mu(\cdot))\in\frak A_c(\Bbb R,X)$ и система


$(1.1)$ \ $\varepsilon,\vartheta,\eta$-РЛУ в малом на полутраекторию


$\orb(\wh x;[\tau_0,\infty))$. Тогда для каждой пары


$(\frak x(\cdot),\frak m(\cdot))


\in\cl(\orb_g(q(\cdot);[\tau_0,\infty)))$ при любом фиксированном


$\varepsilon_1\in(0,\varepsilon)$


система $(1.1)$ \ $\varepsilon_1,\vartheta,\eta$-РЛУ на


$\orb(\frak x;[\tau_0,\infty))$.


\end{lems}








\begin{pf}


Так как $\frak p(\cdot)\in\cl(\orb_g(\wh q(\cdot);[\tau_0,\infty)))$,


то найдется такая последовательность


$\{\tau_j\}_{j=1}^{\infty}\subset[\tau_0,\infty)$, что


$\lim\limits_{j\to\infty}\varrho_{c,w}((\frak p(\cdot),


q_{\tau_j}(\cdot))){=}0$.


Фиксируем $\varepsilon_1\in(0,\varepsilon)$, $\tau\ge\tau_0$ и


$x_0\in O_{\varepsilon_1}[\frak x(\tau)]$. Из предыдущего


предельного равенства и леммы 3.7 (см. также замечание 3.2) получаем


\begin{equation}


\lim_{j\to\infty}(\max_{t\in[\tau,\tau+\vartheta]}


|\frak x(t)-\wh x_{\tau_j}(t)|+


\|\frak m-\wh\mu_{\tau_j}\|_{w,[\tau,\tau+\vartheta]})=0.


\end{equation}


Следовательно, существует $\wh j\in\Bbb N$, начиная с которого


$x_0\in O_{\varepsilon}[\wh x_{\tau_j}(\tau)]$. Далее, т.\,к. система (1.1)


$\varepsilon,\eta,\vartheta$-РЛУ на


$\orb(\wh x;[\tau_0,\infty))$, то для каждого $j\ge\wh j$


найдется управление $\nu_j\in\cal M_{\tau,\vartheta}$, удовлетворяющее


неравенству


\begin{equation}


\|\wh\mu_{\tau_j}-\nu_j\|_{w,[\tau,\tau+\vartheta]}


\le\eta|\wh x_{\tau_j}(\tau)-x_0|,


\end{equation}


при котором система (1.1) имеет решение (см. (2.2))


$y_j(t)=x_0+\int\limits_\tau^t\langle\nu_j(s),f(y_j(s),u)\rangle ds\in K$,


$t\in[\tau,\tau+\vartheta]$, такое, что


\begin{equation}


y_{j}(\tau)=x_0,\ \ y_j(\tau+\vartheta)=x_{\tau_{j}}(\tau+\vartheta).


\end{equation}





Рассмотрим последовательность


$\{y_j(\cdot),\nu_j(\cdot)\}_{j=1}^{\infty}


\subset\frak A_c([\tau,\tau+\vartheta];K)$.


По лемме 1.1 из нее можно выделить подпоследовательность


$\{y_{j_i}(\cdot),\nu_{j_i}(\cdot)\}_{i=1}^{\infty}$, сходящуюся


при $i\to\infty$ в пространстве


$C[\tau,\tau+\vartheta]\times\cal M_{\tau,\vartheta}$ к некоторому


процессу $(y(\cdot),\nu(\cdot))\subset\frak A_c([\tau,\tau+\vartheta];K)$.


Принимая во внимание (4.2)--(4.4), получаем, что для каждой точки


$x_0\in O_{\varepsilon_1}[\frak x(\tau)]$ существует


$\nu\in\cal M_{\tau,\vartheta}$, удовлетворяющее неравенству


$\|\frak m-\nu\|_{w,[\tau,\tau+\vartheta]}


\le\eta|\frak x(\tau)-x_0|$,


при котором система (1.1) имеет решение


$y(t)\in K$, $t\in[\tau,\tau+\vartheta]$, и


такое, что $y(\tau)=x_0$, $y(\tau+\vartheta)={\frak x}(\tau+\vartheta)$.


\end{pf}





Аналогично доказывается








\begin{lems}


Пусть $q(\cdot)\doteq(\wh x(\cdot),\wh\mu(\cdot))


\in\frak A_c(\Bbb R,X)$ и система


$(1.1)$ \ $\varepsilon,\vartheta,\eta$-РЛД в малом с полутраектории


$\orb(\wh x;[\tau_0,\infty))$. Тогда


для каждой пары $(\frak x(\cdot),\frak m(\cdot))


\in\cl(\orb_g(q(\cdot);[\tau_0,\infty)))$ при любом фиксированном


$\varepsilon_1\in(0,\varepsilon)$


система (1.1) \ $\varepsilon_1,\vartheta,\eta$-РЛД


с


$\orb(\frak x;[\tau_0,\infty))$.


\end{lems}








\begin{lems}


Пусть $(\wh x(\cdot),\wh\mu(\cdot))\in O\cal P(\Bbb R)$.


Тогда $(\wh x_{\tau}(\cdot),


\wh\mu_{\tau}(\cdot))\in O\cal P(\Bbb R)$ при каждом $\tau\in\Bbb R$.


\end{lems}








\begin{pf}


Допустим, что найдутся точки $t_0<t_1$ и процесс


$(y(\cdot),\nu(\cdot))\in\frak A[t_0,t_1]$ такие, что


$y(t_0)=x_{\tau}(t_0)$, $y(t_1)=x_{\tau}(t_1)$


и


$\frak T(y(\cdot),\nu(\cdot);t_0,t_1)<


\frak T(\wh x_{\tau}(\cdot),\wh\mu_{\tau}(\cdot);t_0,t_1)$.


Теперь, если рассмотреть $\wt y(t)\doteq y(t-\tau)$,


$\wt \nu(t)\doteq\nu(t-\tau)$, $t\in[t_0+\tau,t_1+\tau]$, то получим


$\frak T(\wt y(\cdot),\wt\nu(\cdot);t_0+\tau$, $t_1+\tau)=


%%%%% join in Eng.                          ^^^^^


\frak T(y(\cdot),\nu(\cdot);t_0,t_1)


<\frak T(\wh x(\cdot),\wh\mu(\cdot);t_0+\tau,t_1+\tau)$,


а это противоречит (см. (1.4)) тому, что $(\wh x(\cdot),\wh\mu(\cdot))


\in O\cal P(\Bbb R)$.


\end{pf}





\addtocounter{section}{1}


\setcounter{equation}{0}


\setcounter{lems}{0}


\setcounter{defns}{0}


\setcounter{notes}{0}


\setcounter{thms}{0}


\setcounter{cors}{0}


\setcounter{exams}{0}





{\bf 5.} Приведем основные утверждения работы.





\begin{thms}


Пусть функция $f:G\times\frak U\to\Bbb R^n$ удовлетворяет условию~$1)$ и


на процессе $q(\cdot)\doteq(\wh x(\cdot),\wh\mu(\cdot))


\in\frak A_{c}(\Bbb R;X)$


выполнено условие~$2)$. Тогда, если система $(2.1)$ РЛУ на $[\tau_0,\infty)$


и $q(\cdot)\in O\cal P[T_1,\infty)$ $(T_1\ge\tau_0)$, то


найдется такое $T_0>T_1$, что всякий процесс $\frak


p(\cdot)\doteq(\frak x(\cdot),\frak m(\cdot))$ из 


$\cl(\orb_g(q(\cdot);[T_1,\infty)))$ принадлежит $O\cal P[T_0,\infty)$.


\end{thms}








{\bf Доказательство.} 


По теореме 2.1 найдутся такие $\varepsilon,\vartheta,\eta>0$, что


система (1.1) будет $\varepsilon,\vartheta,\eta$-РЛУ в малом на


$\orb(\wh x;[\tau_0,\infty))$,


$\varepsilon,\vartheta,\eta$-РЛД в малом с


$\orb(\wh x;[\tau_0,\infty))$ и будет, конечно, обладать этими же 


свойствами относительно $\orb_+(\wh x;[T_1,\infty))$.


Поэтому по леммам 4.1 и 4.2 для каждого процесса


$\frak p(\cdot)\in\cl(\orb_g(q(\cdot);[T_1,\infty)))$


система (1.1) также будет $\varepsilon_1,\vartheta,\eta$-РЛУ


на $\orb(\frak x;[T_1,\infty))$


и $\varepsilon_1,\vartheta,\eta$-РЛД


с $\orb(\frak x;[T_1,\infty))$ при любом


фиксированном $\varepsilon_1\in(0,\varepsilon)$. Кроме того, найдется


такая последовательность $\{\tau_j\}_{j=1}^{\infty}\subset[T_1,\infty)$,


что


\begin{equation}


\lim_{j\to\infty}\varrho_{c,w}(q_{\tau_j}(\cdot),\frak p(\cdot))=0.


\end{equation}


При этом, поскольку $q(\cdot)\in O\cal P[T_1,\infty)$, то


$q_{\tau_j}(\cdot)\in


O\cal P[T_1,\infty)$ по лемме 4.3 при каждом $j{\in}\Bbb N$, следовательно,


$q_{\tau_j}(\cdot)|_{\Bbb T}\in


O\cal P(\Bbb T)$ для каждого отрезка $\Bbb T\subset[T_1,\infty)$.


Поэтому, зафиксировав любую систему отрезков


$\Bbb T_1\subset\Bbb T_2\subset\dotsb$, исчерпывающих


$[T_1,\infty)$, в силу (5.1) и леммы 3.7 получаем, что на каждом


отрезке $\Bbb T_i$, $i\in\Bbb N$, процесс $\frak p(\cdot)$


является пределом в пространстве


$C(\Bbb T_i)\times\cal M_{\Bbb T_i}$ последовательности


$\{q_{\tau_j}(\cdot)|_{\Bbb T_i}\}_{j=1}^{\infty}


\subset O\cal P(\Bbb T_i)$. Отсюда, учитывая, что система (1.1)


является $\varepsilon_1,\vartheta,\eta$-РЛУ на


$\orb(\frak x;[T_1,\infty))$


и $\varepsilon_1,\vartheta,\eta$-РЛД


с $\orb(\frak x;[T_1,\infty))$, так же, как в теореме~2.2, получаем, что


$\frak p(\cdot)\in O\cal P[T_0,\infty)$ при $T_0\doteq T_1+\vartheta$.





Далее, в силу свойств динамической системы сдвигов, определенной на


$(C_b(\Bbb R,\Bbb R^n),\varrho_c)$ (см. [16], c.\,61), и


теоремы 3.4 при $\frak Y\doteq\frm(\frak U)$, принимая во


внимание определение метрики $\varrho_{c,w}$ на


$C_b(\Bbb R,\Bbb R^n)\times\cal M_{\Bbb R}$ (см. (4.1)),


получаем, что для каждого процесса


$q(\cdot)\in\frak A_c(\Bbb R;X)$, в котором $\mu(\cdot)$


является $d_{w}$-непрерывным, множество


$\cl(\orb_g^+(q(\cdot))\in\comp(\frak P))$.


Поэтому по теореме существования минимальных множеств ([15], с.\,307)


множество омега-предельных точек движения $t\mapsto g^t(q(\cdot))$,


названных в [8], [21] магистральными процессами системы (1.1), а в


[11] --- асимптотическими достижимыми процессами этой системы, содержит


минимальное компактное множество $\cal P(q(\cdot))$. Поскольку по


теореме Биркгофа всякое движение динамической системы, начинающееся в


минимальном компактном множестве, рекуррентно, то из теоремы 5.1 и


леммы 3.4, принимая во внимание, что для каждого


$\frak p(\cdot)\in\cal P(q(\cdot))$ при любом фиксированном


$T>0$ \ $\cl(\orb_g({\frak p}(\cdot)))=


\cl(\orb_g(\frak p(\cdot);[T,\infty)))$, получаем








\begin{cors}


Пусть в условиях теоремы 5.1 \ $\wh\mu(\cdot)$ \ $d_{w}$-непрерывно.


Тогда всякий процесс $\frak p(\cdot)\in\cal P(q(\cdot))$ является 


рекуррентным и принадлежит $O\cal P(\Bbb R)$.


\end{cors}





В $\cal M_{\Bbb R}$ выделим подмножество


$\APM_1$ [9], состоящее из таких измеримых отображений


$\mu:\Bbb R\to\rpm(\frak U)$, что для любой функции


$c\in C(\frak U,\Bbb R)$ отображение


$t\mapsto\langle\mu(t),c(u)\rangle\doteq\int\limits_{\frak U}c(u)\mu(t)(du)$


принадлежит пространству $S(\Bbb R,\Bbb R)


\subset L_1^{\loc}(\Bbb R,\Bbb R)$


п.\,п. по Степанову функций [13], а в


$\frak A_c(\Bbb R)$ (см. раздел~1) выделим подмножество


$\cal A_c$, состоящее из таких пар


$(x(\cdot),\mu(\cdot))$, в которых $x(t)$, $t\in\Bbb R$,


--- п.\,п. по Бору [13] решение системы (1.1), отвечающее


управлению $\mu(\cdot)\in\APM_1$.





В дальнейшем говорим, что процесс $\frak p(\cdot)=


(\frak x(\cdot),\frak m(\cdot))\in\frak P$ устойчив


(по Ляпунову) относительно отрицательных сдвигов в положительном


направлении, если для любого $\varepsilon<0$ найдутся такие


$\delta>0$ и $T>0$, что для каждого $\vartheta<0$, при котором


$\max\limits_{t\in[0,T]}\Big\{|\frak x_\vartheta(t)-\frak


x(t)|+\int\limits_t^{t+1}\rho_w({\frak m}_{\vartheta}(s),\frak


m(s))ds\Big\}<\delta$, при всех $t\ge T$ будет выполняться


неравенство


$$


|\frak x_\vartheta(t)-\frak x(t)|


+\int_t^{t+1}\rho_w({\frak m}_{\vartheta}(s),{\frak m}(s))ds


<\varepsilon.


$$








\begin{thms}


Пусть в условиях теоремы $5.1$


управление $\wh\mu(\cdot)$ в процессе $q(\cdot)\doteq(\wh


x(\cdot),\wh\mu(\cdot))$ является $d_{w}$-непрерывным. Тогда всякий


процесс $\frak p(\cdot)\doteq(\frak x(\cdot),


\frak m(\cdot))$, принадлежащий минимальному подмножеству $\cal P(q(\cdot))$


множества омега-предельных точек движения


$t\mapsto g^t(q(\cdot))$, устойчивый относительно


отрицательных сдвигов в положительном направлении,


принадлежит~$\cal A_c$.


\end{thms}








\begin{pf}


По следствию 5.1


отображение $t\mapsto\frak p(t)=(\frak x(t),\frak m(t))$ рекуррентно. 


Следуя рассуждениям доказательства теоремы Маркова ([15], c.\,416),


принимая во внимание определение метрики $\varrho_{c,w}$ и устойчивости


$\frak p(\cdot)$ относительно отрицательных сдвигов в положительном


направлении, как и в [26], можно показать, что для каждого $\varepsilon>0$


множество (см. (3.6)) $\{\tau\in\Bbb R:\sup\limits_{t\in\Bbb R}


|\frak x_\tau(t)-\frak x(t)|+


d_w(\frak m_\tau,\frak m)\le\varepsilon\}$


относительно плотно. Следовательно, функция $t\mapsto\frak x(t)$ п.\,п. по


Бору, а отображение $\frak m(\cdot)\in S(\Bbb R,\rpm(\frak U))$.


Последнее равносильно [9] тому, что $\frak m(\cdot)$ принадлежит


пространству $\APM_1$   %%%%%  TAKE INTO ACCOUNT THE MEANING OF \APM!


мерозначных п.\,п. отображений. Таким


образом, $\frak p(\cdot)\in\cal A_c$.


\end{pf}








\begin{notes}


В теореме 5.2 приведено достаточное


условие почти периодичности рекуррентного движения в терминах


устойчивости $\frak p(\cdot)$ относительно отрицательных сдвигов


в положительном направлении. Действительно, с учетом свойств динамической


системы сдвигов, определенной на


$(C_b(\Bbb R,\Bbb R^n),\varrho_c)$ [15], [16], и


приведенных в разделе 3 утверждений о ее свойствах для случая


$\frak Y\doteq\rpm(\frak U)$ можно указать для определенной


динамической системы $(\frak P,g^t)$ (см. (5.1)) и другие достаточные


условия, используя достаточные условия почти периодичности движений


заданной динамической системы $({\frak X},f^t)$ (напр., 


[15], [16], [27]).


\end{notes}





\input ivanov2.tex





\vskip 2cm





\post{\it Институт математики и информатики}{\it Поступила}


\post{\it Удмуртского государственного}{06.07.2004}


\post{\it университета}





\label{end}


\end{document}


